\subsection{Benutzerinteraktion über die Webschnittstelle}
Zur Steuerung der WordClock wird ein integrierter Webserver eingesetzt. 
Dieser ermöglicht die Anzeige von Statusinformationen sowie die Anpassung von Parametern wie Helligkeit, Farbe und Betriebszustand in einem \gls{ui}. 
Die Verarbeitung der Anfragen erfolgt ereignisbasiert und ist unabhängig vom zyklischen Hauptprogramm. 
Statusinformationen werden bei Bedarf dynamisch im \gls{json}-Format erzeugt und an den Client übertragen.
Der Ablauf ist in Abbildung~\ref{fig:webserver} dargestellt.

\begin{figure}[H]
\centering
\begin{tikzpicture}[
    node distance=1.1cm,
    startstop/.style={rectangle, rounded corners, minimum width=4cm, draw=black},
    process/.style={rectangle, minimum width=4.8cm, draw=black},
    decision/.style={diamond, aspect=2.3, draw=black},
    io/.style={trapezium, trapezium left angle=70, trapezium right angle=110, minimum width=4.6cm, minimum height=0.8cm, text centered, draw=black},
    arrow/.style={->, thick}
]

\node (start) [startstop] {HTTP-Anfrage};
\node (path) [decision, below=of start] {Anfragepfad?};

\node (html) [process, below left=1.5cm and 1.8cm of path] {HTML-Seite erzeugen};
\node (json) [process, below=of path] {JSON-Status erzeugen};
\node (set) [process, below right=1.5cm and 1.8cm of path] {Parameter setzen};

\node (end) [startstop, below=3cm of path] {Antwort senden};

\draw [arrow] (start) -- (path);
\draw [arrow] (path) -- (html);
\draw [arrow] (path) -- (json);
\draw [arrow] (path) -- (set);

\draw [arrow] (html) |- (end);
\draw [arrow] (json) -- (end);
\draw [arrow] (set) |- (end);

\end{tikzpicture}
\caption[Schematischer Programmablaufplan Webserveranfragen]{Schematischer Programmablaufplan für die Verarbeitung der Webserveranfragen über die \gls{ui} der WordClock.}
\label{fig:webserver}
\end{figure}